\section{Conclusiones}

El objetivo general se cumplió al formular y resolver numéricamente la conducción estacionaria en la placa idealizada y comparar los solucionadores sobre el mismo sistema discretizado. En respuesta al primer objetivo específico, las diferencias centrales condujeron a una matriz dispersa simétrica definida positiva, con condiciones de Dirichlet y una discretización que conservó las potencias prescritas. Para el escenario base del segundo objetivo, la solución directa produjo $T_{\max}=77.7592\ ^\circ$C; con residuo relativo menor que $10^{-8}$, Gauss--Seidel convergió en 2145 iteraciones y Gradiente Conjugado en 127. Respecto del tercer objetivo, el refinamiento redujo las diferencias nodales entre mallas sucesivas y el aumento de potencia de las fuentes elevó el máximo térmico. Para el cuarto, ambos métodos alcanzaron la tolerancia y coincidieron estrechamente con la referencia directa, mientras CG requirió menos iteraciones y presentó menores tiempos medianos en las corridas registradas.

La contribución del estudio es una implementación reproducible que separa el error algebraico del cambio espacial y de las simplificaciones físicas, y que compara los métodos con repeticiones de tiempo y exportaciones trazables. El alcance se limita a una placa homogénea bidimensional en estado estacionario, con fuentes uniformes idealizadas y temperatura de contorno fija; por ello, los resultados no predicen temperaturas de una laptop real ni incluyen convección, radiación, transitorios o materiales multicapa. Como continuación, convendría incorporar propiedades y geometría de una placa documentada, condiciones de disipación medidas y un modelo transitorio, y contrastar las predicciones con datos experimentales antes de atribuirles valor de diseño.
